\documentclass[aspectratio=169,12pt]{beamer}
\usepackage[T1]{fontenc}
\usepackage[utf8]{inputenc}
\usepackage{lmodern}
\usepackage{listings}
\usepackage{pgfpages}
\usepackage{graphicx}
\definecolor{Forest}{HTML}{006747}
\definecolor{Ink}{HTML}{17312B}
\definecolor{CodeBg}{HTML}{F2F6F4}
\setbeamercolor{normal text}{fg=Ink,bg=white}
\setbeamercolor{structure}{fg=Forest}
\setbeamercolor{frametitle}{fg=Forest}
\setbeamercolor{alerted text}{fg=Forest}
\setbeamertemplate{navigation symbols}{}
% DLSU CCS branding. Keep the PNG beside this source when sharing.
\setbeamertemplate{footline}{%
  \leavevmode\hbox to \paperwidth{%
    \hspace{5mm}%
    \raisebox{-1mm}{\includegraphics[height=5mm]{DLSU-CCS-Logo-Clear.png}}%
    \hspace{3mm}\scriptsize DLSU College of Computer Studies%
    \hfill HTML5 Foundations\quad\insertframenumber/\inserttotalframenumber\hspace{5mm}%
  }\vspace{2mm}%
}
\titlegraphic{\includegraphics[height=13mm]{DLSU-CCS-Logo-Clear.png}}
\institute{De La Salle University\\College of Computer Studies}
\setbeamersize{text margin left=9mm,text margin right=9mm}
\setbeamerfont{frametitle}{size=\LARGE,series=\bfseries}
\setbeameroption{hide notes}
% For presenter notes, replace the line above with:
% \setbeameroption{show notes on second screen=right}
% For a printable notes document, use: \setbeameroption{show only notes}
\lstset{language=HTML,basicstyle=\ttfamily\normalsize,
 keywordstyle=\color{Forest}\bfseries,commentstyle=\color{gray},
 stringstyle=\color{blue!55!black},backgroundcolor=\color{CodeBg},
 showstringspaces=false,columns=fullflexible,keepspaces=true,
 breaklines=true,tabsize=4,frame=none,aboveskip=6pt,belowskip=6pt}
\title{HTML5 Foundations}
\subtitle{Building a student portfolio}
\author{Introductory college lesson}
\date{}
\begin{document}
% Slide 01: CORE: 35 seconds
\begin{frame}[plain]
\relax
{\usebeamerfont{title}\color{Forest}\LARGE\bfseries HTML5 Foundations\par}
\medskip
{\large Building a student portfolio\par}
\medskip
\includegraphics[height=13mm]{DLSU-CCS-Logo-Clear.png}
\medskip
\textbf{By the end, you can:}
\begin{itemize}
\item Write a valid HTML document with meaningful content.
\item Organize a portfolio with semantic sections.
\end{itemize}

\note{CORE: 35 seconds. Introduce Maya, a fictional student who needs one portfolio page. The end product is a readable page, using default browser styles. Explain that we will add CSS next week.}
\end{frame}
% Slide 02: CORE: 25 seconds
\begin{frame}[fragile]{Generative AI as a learning assistant}
\large Generative AI can assist your learning. Understand every line you use, verify its output, and be able to explain and modify the HTML yourself.
\par\medskip\normalsize
\textbf{A useful check:} can you predict what the browser will display?
\note{CORE: 25 seconds. Ask students to use AI for explanations and feedback, then test the result themselves. Ask what an unfamiliar tag means before copying it. Do not introduce an AI account or tool setup.}
\end{frame}
% Slide 03: CORE: 45 seconds
\begin{frame}[fragile]{A webpage in a browser}
\begin{itemize}
\item An HTML file contains text that describes page content.
\item The browser interprets the tags and displays their content; the tags themselves are not shown.
\item Changing the saved file and refreshing changes what you see.
\end{itemize}

\note{CORE: 45 seconds. Use the familiar browser tab as the starting point. Explain that local files work for today, while a web server can deliver the same kinds of files online. Avoid suggesting that browsers only render HTML. Source: \url{https://developer.mozilla.org/en-US/docs/Learn_web_development/Core/Structuring_content/Basic_HTML_syntax}.}
\end{frame}
% Slide 04: CORE: 40 seconds
\begin{frame}[fragile]{HTML describes content}
\begin{lstlisting}
<h1>Maya Santos</h1>
<p>Computing student and curious learner.</p>
\end{lstlisting}
\par\medskip \textbf{Quick discussion:} HTML means HyperText Markup Language.
It describes webpage structure: \texttt{h1} marks a main heading and \texttt{p} a paragraph.
\par\smallskip {\small Introductory reference: \href{https://www.w3schools.com/html/html_intro.asp}{W3Schools: HTML Introduction}.}

\note{CORE: 40 seconds. Read h1 as a main heading and p as a paragraph. Ask which text looks like a heading. Describe markup as adding meaningful labels, rather than writing calculations or decision logic. Source: \url{https://developer.mozilla.org/en-US/docs/Learn_web_development/Core/Structuring_content/Basic_HTML_syntax}.}
\end{frame}
% Slide 05: CORE: 40 seconds
\begin{frame}[fragile]{Why HTML5 before frameworks?}
\begin{itemize}
\item \textbf{Bootstrap builds on HTML:} its ready-made styles and components use HTML elements.
\item \textbf{Choose meaningful structure:} headings, links, and page regions still need the right tags.
\item \textbf{Debug with understanding:} recognize broken links, nesting errors, and missing attributes.
\end{itemize}
\par\medskip Learn the structure first, then use a framework to build on it.
\note{CORE: 40 seconds. Use Bootstrap as the concrete example of a frontend UI framework/toolkit. It supplies CSS and JavaScript features, but developers still write the HTML structure. No framework setup or code is needed today. This is a recommended learning order, not a requirement to master all HTML before beginning a framework. Connect this back to the portfolio: its headings and section links still need meaningful markup after styling. Source: \url{https://getbootstrap.com/docs/5.3/getting-started/introduction/}.}
\end{frame}
% Slide 06: OPTIONAL: about 60 seconds; excluded from the 20-minute core route
\begin{frame}[fragile]{HTML, CSS, and JavaScript}
\begin{itemize}
\item \textbf{HTML:} content and structure, such as a heading.
\item \textbf{CSS:} appearance, such as color and spacing.
\item \textbf{JavaScript:} behavior, such as responding to an action.
\end{itemize}
\par\medskip Today we focus on HTML. Browser defaults supply the appearance.

\note{OPTIONAL: about 60 seconds; excluded from the 20-minute core route. Optional context. Use one heading as the example for all three roles. HTML can provide built-in interactions such as links; JavaScript is not required for every interaction. Do not teach CSS rules or JavaScript code.}
\end{frame}
% Slide 07: CORE: 40 seconds
\begin{frame}[fragile]{Our portfolio destination}
\begin{itemize}
\item \textbf{Introduction:} Maya\textquotesingle s name and a short description.
\item \textbf{Main content:} About, Skills, and Projects.
\item \textbf{Navigation:} links to those sections.
\item \textbf{Footer:} a short closing note.
\end{itemize}

\note{CORE: 40 seconds. Briefly show the completed local page already open before class, then return to slides. Emphasize that we will build the structure ourselves, not copy a visual template. The supplied college image is a logo, not a portrait.}
\end{frame}
% Slide 08: CORE: 50 seconds
\begin{frame}[fragile]{An element and its tags}
\begin{lstlisting}
<p>I am learning HTML.</p>
\end{lstlisting}
\begin{itemize}
\item \texttt{<p>} starts the paragraph.
\item \texttt{I am learning HTML.} is its content.
\item \texttt{</p>} ends the paragraph.
\end{itemize}
\par\smallskip \textbf{Quick discussion:} The tags and content together form
one element. Change the text, save, and refresh to see the result.


\note{CORE: 50 seconds. An element includes the tags and content. Point at the slash in the closing tag. The next supporting slide explains the exception for void elements. Source: \url{https://developer.mozilla.org/en-US/docs/Learn_web_development/Core/Structuring_content/Basic_HTML_syntax}.}
\end{frame}
% Slide 09: OPTIONAL: about 60 seconds; excluded from the 20-minute core route
\begin{frame}[fragile]{Elements without closing tags}
\begin{lstlisting}
<img src="assets/college-logo.png"
     alt="DLSU and CCS emblems">
\end{lstlisting}
\par\medskip \texttt{img} is a void element. It has no closing tag or nested content.

\note{OPTIONAL: about 60 seconds; excluded from the 20-minute core route. Optional clarification. The attributes will be taught next, and images later. Explain that a closing img tag should not be added. meta is another void element we will meet in the document head. A self-closing slash is not required in HTML. Source: \url{https://developer.mozilla.org/en-US/docs/Web/HTML/Reference/Elements/img}.}
\end{frame}
% Slide 10: CORE: 45 seconds
\begin{frame}[fragile]{Nesting creates a hierarchy}
\begin{lstlisting}
<section>
  <h2>About me</h2>
  <p>I am Maya, a computing student.</p>
</section>
\end{lstlisting}
\par\medskip Children sit inside a parent. Close the inner element before the outer one.

\note{CORE: 45 seconds. Section and h2 roles are coming later; here focus on structure. Point out indentation as a reading aid. Indentation makes nesting easier to read; the tags determine the structure. Ask which element contains the paragraph. Source: \url{https://developer.mozilla.org/en-US/docs/Learn_web_development/Core/Structuring_content/Basic_HTML_syntax}.}
\end{frame}
% Slide 11: CORE: 40 seconds
\begin{frame}[fragile]{Attributes add information}
\begin{lstlisting}
<a href="#projects">My projects</a>
\end{lstlisting}
\begin{itemize}
\item \texttt{href} is the attribute name.
\item \texttt{\#projects} is its quoted value.
\item Attributes appear in the opening tag.
\end{itemize}

\note{CORE: 40 seconds. Explain that an attribute supplies information about an element. The fragment destination will make sense after IDs. Ask which part is visible text. Answer My projects, not the href value. Source: \url{https://developer.mozilla.org/en-US/docs/Learn_web_development/Core/Structuring_content/Basic_HTML_syntax}.}
\end{frame}
% Slide 12: CORE: 80 seconds
\begin{frame}[fragile]{The HTML5 document skeleton}
\begin{lstlisting}
<!DOCTYPE html>
<html lang="en">
<head>
  <title>Maya's portfolio</title>
</head>
<body>
  <h1>Maya Santos</h1>
</body>
</html>
\end{lstlisting}
\par\smallskip {\small \textbf{Quick discussion:} The doctype selects standards mode;
\texttt{html} is the root and \texttt{lang} identifies its language.
\texttt{head} holds metadata; \texttt{body} holds displayed content.
The title appears in the tab; \texttt{h1} appears on the page.}

\note{CORE: 80 seconds. Allow 80 seconds on the core route. Explain doctype as standards mode, html as root, lang as the document language, head as metadata and body as displayed content. The title names the browser tab, while h1 names the page content. Full files include charset and viewport metadata shown on the next supporting slide. HTML5 is common teaching terminology; HTML now develops as a Living Standard. Source: \url{https://html.spec.whatwg.org/multipage/}.}
\end{frame}
% Slide 13: OPTIONAL: about 60 seconds; excluded from the 20-minute core route
\begin{frame}[fragile]{Metadata in the head}
\begin{lstlisting}
<meta charset="UTF-8">
<meta name="viewport"
      content="width=device-width, initial-scale=1.0">
<title>Maya's portfolio</title>
\end{lstlisting}
\begin{itemize}
\item UTF-8 sets the document character encoding.
\item The viewport setting supports mobile-sized layouts.
\item The title helps identify the browser tab.
\end{itemize}

\note{OPTIONAL: about 60 seconds; excluded from the 20-minute core route. Optional detail. These lines live inside head in all complete examples. Viewport metadata alone does not make a layout responsive. It prepares a normal mobile viewport for later CSS. meta has no closing tag. Source: \url{https://developer.mozilla.org/en-US/docs/Learn_web_development/Core/Structuring_content/Basic_HTML_syntax}.}
\end{frame}
% Slide 14: CORE: 55 seconds
\begin{frame}[fragile]{Headings describe the outline}
\begin{lstlisting}
<h1>Maya Santos</h1>
<h2>Projects</h2>
<h3>Campus study guide</h3>
\end{lstlisting}
\par\medskip \textbf{Quick discussion:} HTML has six heading levels,
\texttt{h1}--\texttt{h6}. Here the page title contains a Projects topic,
which contains a project title. Choose levels for hierarchy, rather than text size.

\note{CORE: 55 seconds. Use one main h1 in this lesson. Projects is a section, and the project title nests below it. Ask why choosing h3 just for a smaller font is unhelpful. CSS will control size next week. Source: \url{https://developer.mozilla.org/en-US/docs/Learn_web_development/Core/Structuring_content/Basic_HTML_syntax}.}
\end{frame}
% Slide 15: OPTIONAL: about 60 seconds; excluded from the 20-minute core route
\begin{frame}[fragile]{Paragraphs organize text}
\begin{lstlisting}
<p>I am Maya, a computing student.</p>
<p>I enjoy making useful webpages.</p>
\end{lstlisting}
\par\medskip \textbf{Quick discussion:} Each \texttt{p} groups a related idea.
Source-code line breaks do not create new paragraphs; use another \texttt{p}.
\par\smallskip \textbf{Try it:} Add a paragraph about another interest, then refresh.

\note{OPTIONAL: about 60 seconds; excluded from the 20-minute core route. Optional practice. Explain ordinary whitespace generally collapses in HTML rendering. Use p elements instead of repeated br tags to divide paragraphs. HTML source can wrap across lines for readability. Source: \url{https://developer.mozilla.org/en-US/docs/Learn_web_development/Core/Structuring_content/Basic_HTML_syntax}.}
\end{frame}
% Slide 16: OPTIONAL: about 60 seconds; excluded from the 20-minute core route
\begin{frame}[fragile]{Importance and emphasis}
\begin{lstlisting}
<p>Please <strong>save your work</strong>.</p>
<p>I am <em>still learning</em> HTML.</p>
\end{lstlisting}
\begin{itemize}
\item \texttt{strong} marks importance.
\item \texttt{em} marks stress emphasis.
\end{itemize}
\par\smallskip \textbf{Quick discussion:} These elements add meaning inside
a paragraph. Close each inner element before closing \texttt{p}.


\note{OPTIONAL: about 60 seconds; excluded from the 20-minute core route. Optional enrichment. Browsers usually render these as bold and italic, but the meaning matters. Avoid teaching strong and em as mere visual formatting commands. Source: \url{https://developer.mozilla.org/en-US/docs/Learn_web_development/Core/Structuring_content/Basic_HTML_syntax}.}
\end{frame}
% Slide 17: CORE: 50 seconds
\begin{frame}[fragile]{Lists organize related items}
\begin{lstlisting}
<ul>
  <li>Writing HTML</li>
  <li>Checking links</li>
</ul>
\end{lstlisting}
\par\medskip \textbf{Quick discussion:} Use \texttt{ul} when order does not
matter, as with skills. Use \texttt{ol} for a sequence, such as edit, save,
and refresh. Both contain \texttt{li} items.
\par\smallskip \textbf{Try it:} Add another skill inside this list.

\note{CORE: 50 seconds. Use the portfolio skills as the example. Ask when order matters: editing, saving and refreshing are ordered steps. The independent list example demonstrates both ul and ol. Source: \url{https://developer.mozilla.org/en-US/docs/Learn_web_development/Core/Structuring_content/Basic_HTML_syntax}.}
\end{frame}
% Slide 18: CORE: 40 seconds
\begin{frame}[fragile]{Links connect destinations}
\begin{lstlisting}
<a href="index.html">Read my portfolio</a>
\end{lstlisting}
\par\medskip \textbf{Quick discussion:} \texttt{a} creates a link;
\texttt{href} gives its destination. The text between the tags is the visible label.
A destination can be a local file, a full URL, or a same-page fragment.

\note{CORE: 40 seconds. This relative link points to a file in the same folder. A full URL can point to another website. Tell students that linking to a page is not the same as embedding it. Avoid opening external sites during this offline lesson. Source: \url{https://developer.mozilla.org/en-US/docs/Learn_web_development/Core/Structuring_content/Basic_HTML_syntax}.}
\end{frame}
% Slide 19: CORE: 50 seconds
\begin{frame}[fragile]{Navigation targets an ID}
\begin{lstlisting}
<a href="#projects">Projects</a>
<section id="projects">
  <h2>Projects</h2>
</section>
\end{lstlisting}
\par\medskip \textbf{Quick discussion:} \texttt{href="\#projects"} targets
\texttt{id="projects"}. Include \# in the link, but not in the ID.
Each ID must be unique.
\par\smallskip \textbf{Try it:} Follow the link with Tab and Enter.

\note{CORE: 50 seconds. Trace projects in both places. ID values should be unique; the href includes the hash while the id value does not. This is the exact pattern used in the assignment. Navigation can work even without CSS or JavaScript. Source: \url{https://developer.mozilla.org/en-US/docs/Learn_web_development/Core/Structuring_content/Basic_HTML_syntax}.}
\end{frame}
% Slide 20: CORE: 50 seconds
\begin{frame}[fragile]{Images need a source and alternative text}
\begin{lstlisting}
<img src="assets/college-logo.png"
     width="240" height="129"
     alt="DLSU and CCS emblems">
\end{lstlisting}
\begin{itemize}
\item \texttt{src} starts from this HTML file and enters \texttt{assets}.
\item \texttt{alt} provides a text alternative for this meaningful image.
\end{itemize}
\par\smallskip \textbf{Quick discussion:} Dimensions are in pixels.
\texttt{img} is a void element: it has no content or closing tag.


\note{CORE: 50 seconds. The local image is supplied and makes the folder work offline. Explain that decorative images use empty alt, but this meaningful institutional mark has descriptive text. The example files use the fuller official names. Image dimensions preserve the image aspect ratio. Source: \url{https://developer.mozilla.org/en-US/docs/Web/HTML/Reference/Elements/img}.}
\end{frame}
% Slide 21: CORE: 45 seconds
\begin{frame}[fragile]{Semantic HTML names the role}
\begin{itemize}
\item Choose elements that explain what the content does.
\item \texttt{header}, \texttt{nav}, \texttt{main}, and \texttt{footer} identify page regions.
\item Meaningful structure helps readers and assistive technology.
\end{itemize}
\par\smallskip \textbf{Quick discussion:} A section groups a topic; an article
is a self-contained item within it. Semantic roles organize meaning;
CSS controls visual layout.


\note{CORE: 45 seconds. We already know how tags nest. Now choose the right tags for each part of the portfolio. Semantics alone do not produce a layout or guarantee accessibility. CSS handles presentation. Source: \url{https://developer.mozilla.org/en-US/docs/Learn_web_development/Core/Structuring_content/Structuring_documents}.}
\end{frame}
% Slide 22: CORE: 35 seconds
\begin{frame}[fragile]{Header: the introduction}
\begin{lstlisting}
<header>
  <h1>Maya Santos</h1>
  <p>Computing student and curious learner.</p>
</header>
\end{lstlisting}
\par\medskip \textbf{Quick discussion:} This page-level \texttt{header}
introduces the portfolio and its owner. It differs from \texttt{head},
which holds document metadata rather than visible introductory content.

\note{CORE: 35 seconds. Use the portfolio introduction students saw at the start. Header can also introduce a section or article; this example is the page-level header. It is different from head, which contains metadata. Source: \url{https://developer.mozilla.org/en-US/docs/Learn_web_development/Core/Structuring_content/Structuring_documents}.}
\end{frame}
% Slide 23: CORE: 35 seconds
\begin{frame}[fragile]{Nav: the major links}
\begin{lstlisting}
<nav aria-label="Portfolio sections">
  <a href="#about">About</a>
  <a href="#projects">Projects</a>
</nav>
\end{lstlisting}
\par\medskip A nav groups a major set of navigation links.

\note{CORE: 35 seconds. Point back to matching IDs. Explain aria-label briefly as a name for this navigation region; it is an attribute, not a replacement for meaningful link text. The complete page uses a list inside nav. Not every group of links needs nav. Source: \url{https://developer.mozilla.org/en-US/docs/Learn_web_development/Core/Structuring_content/Structuring_documents}.}
\end{frame}
% Slide 24: CORE: 35 seconds
\begin{frame}[fragile]{Main: the primary content}
\begin{lstlisting}
<main>
  <h2>About me</h2>
  <p>I am learning to build webpages.</p>
</main>
\end{lstlisting}
\par\medskip Use one main region for this page	extquotesingle s primary content.

\note{CORE: 35 seconds. Show the relationship between the page-level header, nav, main and footer. The full portfolio places About, Skills and Projects sections inside main. Repeated navigation belongs outside this main content. Source: \url{https://developer.mozilla.org/en-US/docs/Learn_web_development/Core/Structuring_content/Structuring_documents}.}
\end{frame}
% Slide 25: CORE: 40 seconds
\begin{frame}[fragile]{Section: a related topic}
\begin{lstlisting}
<section id="skills">
  <h2>Skills</h2>
  <ul><li>Writing HTML</li></ul>
</section>
\end{lstlisting}
\par\medskip A section groups a topic, usually with a descriptive heading.

\note{CORE: 40 seconds. About, Skills and Projects are separate topics inside main. This pattern prepares students to add Interests after class. Use div for a generic grouping without a semantic role when appropriate, rather than making every wrapper a section. Source: \url{https://developer.mozilla.org/en-US/docs/Learn_web_development/Core/Structuring_content/Structuring_documents}.}
\end{frame}
% Slide 26: OPTIONAL: about 60 seconds; excluded from the 20-minute core route
\begin{frame}[fragile]{Article: a self-contained item}
\begin{lstlisting}
<article>
  <h3>Campus study guide</h3>
  <p>A project helping new students.</p>
</article>
\end{lstlisting}
\par\medskip Use article for a self-contained piece, such as one project description.

\note{OPTIONAL: about 60 seconds; excluded from the 20-minute core route. Optional deeper explanation. Projects is the surrounding section; one complete project description can be an article. Not every paragraph or decorative box is an article. Ask whether this description could make sense if shared separately. Source: \url{https://developer.mozilla.org/en-US/docs/Learn_web_development/Core/Structuring_content/Structuring_documents}.}
\end{frame}
% Slide 27: CORE: 25 seconds
\begin{frame}[fragile]{Footer: the closing information}
\begin{lstlisting}
<footer>
  <p>Learning portfolio. Fictional student details.</p>
</footer>
\end{lstlisting}
\par\medskip A footer supplies closing information for its page or section.

\note{CORE: 25 seconds. Use a brief authorship or contextual note. Footer is semantic, not an instruction to pin content to the bottom of the window. Positioning belongs to CSS. Source: \url{https://developer.mozilla.org/en-US/docs/Learn_web_development/Core/Structuring_content/Structuring_documents}.}
\end{frame}
% Slide 28: CORE: 70 seconds
\begin{frame}[fragile]{The portfolio structure together}
\begin{lstlisting}
<header>...</header>
<nav>...</nav>
<main>
  <section id="about">...</section>
  <section id="skills">...</section>
  <section id="projects">...</section>
</main>
<footer>...</footer>
\end{lstlisting}
\par\medskip This is a structure map. Replace the dots with the content already discussed.

\note{CORE: 70 seconds. Allow 70 seconds. Walk through the map from top to bottom. The dots deliberately omit details; the completed example contains real content and the article. Ask students where Interests should go: inside main as a new section, not in head. Source: \url{https://developer.mozilla.org/en-US/docs/Learn_web_development/Core/Structuring_content/Structuring_documents}.}
\end{frame}
% Slide 29: OPTIONAL: about 60 seconds; excluded from the 20-minute core route
\begin{frame}[fragile]{Files that travel together}
\begin{lstlisting}
portfolio/
  index.html
  assets/
    college-logo.png
\end{lstlisting}
\par\medskip A relative image path starts from the HTML file	extquotesingle s location.

\note{OPTIONAL: about 60 seconds; excluded from the 20-minute core route. Optional setup. This listing is a folder map, not HTML code. Keep assets with the HTML when copying or submitting a portfolio. Renaming or moving an image without updating src breaks the path.}
\end{frame}
% Slide 30: OPTIONAL: about 60 seconds; excluded from the 20-minute core route
\begin{frame}[fragile]{Edit, save, and refresh}
\begin{itemize}
\item Open index.html in a plain-text editor.
\item Open the same file in a browser.
\item Save an edit, then refresh the browser.
\item Check that the changed text appears.
\end{itemize}

\note{OPTIONAL: about 60 seconds; excluded from the 20-minute core route. Optional workflow. Demonstrate this during the final live segment rather than spending core time on setup. Ensure the extension is .html rather than .html.txt. A local server is not needed.}
\end{frame}
% Slide 31: OPTIONAL: about 60 seconds; excluded from the 20-minute core route
\begin{frame}[fragile]{Common mistakes to check}
\begin{itemize}
\item Missing or incorrectly nested closing tags.
\item An href fragment that does not match an ID.
\item An image filename or folder path that differs from src.
\item Editing one copy while viewing a different copy.
\end{itemize}

\note{OPTIONAL: about 60 seconds; excluded from the 20-minute core route. Optional troubleshooting. Browsers repair some invalid markup, so a page looking acceptable does not prove the source is correct. Refer to the full examples for comparison. No deliberately broken code is included in learner files. Source: \url{https://developer.mozilla.org/en-US/docs/Learn_web_development/Core/Structuring_content/Basic_HTML_syntax}.}
\end{frame}
% Slide 32: CORE: 50 seconds
\begin{frame}[fragile]{A quick accessibility check}
\begin{itemize}
\item Headings describe the content hierarchy.
\item Link text explains its destination.
\item Meaningful images have useful alternative text.
\item Use Tab and Enter to test navigation links.
\end{itemize}

\note{CORE: 50 seconds. Allow 40 seconds on the core route. Ask students to name one check they can perform without a special tool. This is a starter checklist, not a full accessibility audit. Link it to reviewing AI-generated HTML. Source: \url{https://developer.mozilla.org/en-US/docs/Learn_web_development/Core/Structuring_content/Structuring_documents}.}
\end{frame}
% Slide 33: CORE: 90 seconds
\begin{frame}[fragile]{Live demo: assemble the portfolio}
\begin{itemize}
\item Paste the prepared content into the starter body.
\item Save index.html and refresh the browser.
\item Activate Projects to check its navigation target.
\end{itemize}

\note{CORE: 90 seconds. Exactly 90 seconds. Before class, copy examples/portfolio-demo-body.txt to the clipboard and open examples/portfolio-starter/index.html in both an editor and browser. 0–15 seconds identify the body placeholder; 15–40 replace only its contents; 40–55 save; 55–75 refresh and point to the sections; 75–90 activate Projects and explain the matching ID. Keep document metadata unchanged. Use the completed example as a fallback if an editor fails.}
\end{frame}
% Slide 34: CORE: 60 seconds
\begin{frame}[fragile]{Assignment: add Interests}
\begin{itemize}
\item Add an Interests section inside main.
\item Give it an h2 heading and a list of two interests.
\item Add a navigation link targeting its unique ID.
\end{itemize}
\par\medskip \textbf{Expected:} a new heading and bullet list, reached through Interests.

\note{CORE: 60 seconds. Exactly 60 seconds. Students complete this after class, not inside this one-minute slot. Give exercises/portfolio-starter as the starting folder. 0–25 seconds explain the task; 25–45 state the expected new heading, list and working link; 45–60 recap structure, content and meaningful regions. Ask learners to explain their own code, including AI-assisted lines. A separate solution is provided.}
\end{frame}
% Slide 35: CORE: 30 seconds
\begin{frame}[fragile]{Next week: styling with CSS}
\begin{itemize}
\item Connect a stylesheet and choose selectors.
\item Set colors and typography.
\item Use spacing and the box model.
\end{itemize}
\par\medskip We will style this same portfolio while retaining its HTML structure.

\note{CORE: 30 seconds. Exactly 30 seconds. Show how this lesson becomes the foundation for the next, without adding CSS syntax today. The portfolio will look plain by design until styling is introduced.}
\end{frame}
\end{document}
